\section{A first continuous model: uniform distribution on an interval}

All models constructed so far in this chapter were discrete. Some of them were
finite, such as dice, cards, and urns. One of them was countably infinite: the
waiting time until the first success. These examples are essential, but they do
not yet represent the full scope of probability. Many random quantities are
naturally modelled as points on a continuum.

The purpose of this section is to introduce one continuous model carefully: the
uniform distribution on an interval. We shall not develop the full measure theory
behind continuous probability. That would require a separate course. However, we
must already see the basic change in viewpoint:
\[
\text{in a continuous model, individual points are not the right objects to count.}
\]

Instead of counting elementary outcomes, we measure sizes of sets. For an
interval on the real line, the relevant size is length.

\subsection*{Choosing a point from an interval}

Suppose an experiment produces a real number between \(0\) and \(1\). For
example, imagine an idealized pointer that lands somewhere on a line segment of
unit length. The sample space is
\[
\Omega=[0,1].
\]
An elementary outcome is a single real number
\[
x\in[0,1].
\]

This sample space is not finite and not countably infinite. It is uncountable.
There is no list
\[
x_1,x_2,x_3,\ldots
\]
that contains all real numbers in \([0,1]\). This already makes the model
qualitatively different from coin tosses, dice rolls, card hands, urn selections,
and waiting times.

\begin{remarkbox}
The phrase ``choose a point uniformly from \([0,1]\)'' cannot mean that every
point receives the same positive probability. There are too many points. The
probability is spread continuously over the interval.
\end{remarkbox}

\subsection*{Why points have probability zero}

In a finite uniform model, each elementary outcome has probability
\[
\frac1{|\Omega|}.
\]
This idea cannot be copied directly to \([0,1]\), because \([0,1]\) has
uncountably many points. If each point had the same positive probability, then
the total probability would be far larger than \(1\). If each point had
probability \(0\), then adding point probabilities one by one cannot recover the
probability of an interval.

In the continuous uniform model, every individual point has probability zero:
\[
P(\{x\})=0
\]
for each \(x\in[0,1]\). This does not mean that the experiment cannot produce a
point. It means that probability is not assigned by counting individual points.
The probability of landing in a region is determined by the length of that
region.

For example, the event
\[
A=[0.2,0.5]
\]
has length \(0.3\), so in the uniform model on \([0,1]\) we assign
\[
P(A)=0.3.
\]
The event contains infinitely many points, but its probability is controlled by
how much of the interval it occupies.

\subsection*{Intervals as events}

The simplest events in a continuous interval model are intervals. In the uniform
model on \([0,1]\), the probability of an interval is its length. Thus
\[
P([c,d])=d-c
\]
whenever
\[
0\leq c\leq d\leq 1.
\]
The same value is assigned to
\[
(c,d),\qquad [c,d),\qquad (c,d].
\]
The endpoints do not change the probability, because individual points have
probability zero.

\begin{examplebox}
Let \(X\) be uniformly distributed on \([0,1]\). Then
\[
P(0.25\leq X\leq 0.75)=0.75-0.25=0.5.
\]
Also,
\[
P(0.25<X<0.75)=0.5.
\]
The difference between the closed and open interval consists only of endpoints,
and endpoints have probability zero in this continuous model.
\end{examplebox}

This is a major difference from discrete probability. In a discrete model,
removing one elementary outcome may change the probability. In a continuous
uniform model, removing finitely many points from an interval does not change
its probability.

\subsection*{Uniform distribution on \([a,b]\)}

The same idea works on any finite interval
\[
[a,b],\qquad a<b.
\]
The sample space is
\[
\Omega=[a,b].
\]
The total length of the sample space is
\[
b-a.
\]
If the model is uniform, then probability is proportional to length. For an
interval \([c,d]\subseteq[a,b]\), we define
\[
P([c,d])=\frac{d-c}{b-a}.
\]

The denominator is the length of the whole sample space. The numerator is the
length of the event. Thus the formula is the continuous analogue of the uniform
finite formula
\[
P(A)=\frac{|A|}{|\Omega|}.
\]
The symbol \(|A|\) counted outcomes in a finite space. In the interval model,
length replaces counting.

\begin{definitionbox}
A random point \(X\) has the \textbf{uniform distribution on \([a,b]\)} if
probabilities of intervals are proportional to their lengths:
\[
P(c\leq X\leq d)=\frac{d-c}{b-a}
\]
for every \([c,d]\subseteq[a,b]\).
\end{definitionbox}

\subsection*{Examples}

Let \(X\) be uniformly distributed on \([2,8]\). The whole interval has length
\[
8-2=6.
\]
The probability that \(X\) lies between \(3\) and \(5\) is
\[
P(3\leq X\leq 5)=\frac{5-3}{8-2}=\frac{2}{6}=\frac13.
\]

The probability that \(X\) is at least \(6\) is
\[
P(X\geq 6)=P(6\leq X\leq 8)=\frac{8-6}{8-2}=\frac13.
\]

The probability that \(X\) is exactly \(4\) is
\[
P(X=4)=0.
\]
This last statement is not a contradiction. The experiment produces some exact
value, but no single prescribed value occupies positive length.

\subsection*{Unions of intervals}

Events need not be single intervals. Suppose \(X\) is uniform on \([0,1]\), and
consider the event
\[
A=[0,0.2]\cup[0.7,1].
\]
The two intervals are disjoint. Their lengths are \(0.2\) and \(0.3\). Hence
\[
P(A)=0.2+0.3=0.5.
\]
This is the same additivity principle used in discrete probability: disjoint
parts have probabilities that add. The difference is that the probabilities are
computed from lengths rather than from counts of elementary outcomes.

\subsection*{What are the events in a continuous model?}

For a finite or countable sample space, it is often harmless at this level to
think of every subset of \(\Omega\) as an event. For continuous sample spaces,
this becomes too optimistic. The interval \([0,1]\) has extremely complicated
subsets, and not every subset can be assigned a length in a way that is
compatible with the usual rules of probability.

Therefore, a continuous probability model is not just
\[
\Omega=[0,1].
\]
It also requires a specified collection of admissible events. For the interval
model, the standard choice is the collection of Borel sets, denoted informally by
\[
\mathcal{B}([0,1]).
\]
These are the sets obtained from intervals by the operations that probability
needs: complements, countable unions, and countable intersections.

\begin{definitionbox}
In a continuous model, the event space is usually not taken to be the set of all
subsets of \(\Omega\). Instead, one specifies a suitable collection of events,
usually a \textbf{sigma-algebra}. On the real line, the standard event collection
is the Borel sigma-algebra generated by open intervals.
\end{definitionbox}

For this course, the technical construction is less important than the modelling
message. Intervals, finite unions of intervals, countable unions of intervals,
and sets obtained from them by complements are legitimate events. These are more
than enough for the elementary continuous examples we need now.

\begin{remarkbox}
The word ``Borel'' is a signal that continuous probability requires more care
than discrete probability. We introduce it here so that the model is honest:
continuous probability spaces need a sample space, an event collection, and a
probability assignment.
\end{remarkbox}

\subsection*{The structure of the continuous uniform model}

The continuous uniform model on \([a,b]\) has three pieces:
\[
\Omega=[a,b],
\]
\[
\mathcal{F}=\mathcal{B}([a,b]),
\]
and
\[
P(A)=\frac{\text{length of }A}{b-a}
\]
for events \(A\) whose length is defined in the standard sense.

At this stage, we mainly use this rule for intervals and simple unions of
intervals. The notation \(\mathcal{F}\) reminds us that events form a selected
collection of subsets of \(\Omega\), not necessarily all subsets.

\subsection*{What this continuous model teaches}

The uniform interval model teaches several lessons that are invisible in purely
discrete examples.

First, a sample space may be uncountable. The interval \([0,1]\) contains more
points than any sequence can list.

Second, probability need not be built by assigning positive mass to individual
elementary outcomes. In the uniform continuous model,
\[
P(\{x\})=0
\]
for every point \(x\), while intervals can have positive probability.

Third, the probability of an event is measured by geometric size, not by the
number of points. For the uniform distribution on \([a,b]\), interval
probabilities are relative lengths.

Fourth, the event collection matters. In continuous probability, we must be more
careful about which subsets of the sample space are allowed as events. The
standard answer on the real line is to use Borel sets.

\begin{theorembox}
The continuous uniform model replaces counting by length:
\[
\text{finite uniform model: } P(A)=\frac{|A|}{|\Omega|},
\]
\[
\text{continuous uniform interval model: } P(A)=\frac{\text{length of }A}{\text{length of }\Omega}.
\]
The idea is analogous, but the mathematics is different because the sample space
is uncountable.
\end{theorembox}
